\begin{enumerate}
\item[\textbf{8.}] [[curso=Geometria]] [[tema=angulos]] Indique el valor de verdad de las siguientes proposiciones:£I. El postulado de la medida de un ángulo tiene demostración matemática.£II. El complemento y el suplemento de un ángulo forman un par lineal.£III. El postulado de la construcción de un ángulo está relacionado con un transportador.£IV. El postulado de la adición de ángulos tiene demostración matemática.£A)FVVFæB)FFVFæC)FFFF£D)FVFVææE)FFVF£ [[clave=-]] [[Estado=sin_revisar]]
\item[\textbf{16.}] [[curso=Geometria]] [[tema=conjuntos convexos]] Indique el valor de verdad de siguientes proposiciones:£I. Una recta $\mathcal{L}$, contenida en un plano, determina los semiplanos $S_1$ y $S_2$, luego $S_1\cap S_2=\mathcal{L}$.£II. Sea $T$ una región triangular y $O$ el circuncentro del triángulo, entonces $T-\{O\}$ es siempre un conjunto no convexo.£III. Ningún punto es un conjunto no convexo.£IV. El interior de un ángulo es un conjunto no convexo.£A)FFFFæB)FFVFæC)FFVV£D)FVVVææE)VVVV£ [[clave=-]] [[Estado=sin_revisar]]
\item[\textbf{17.}] [[curso=Geometria]] [[tema=conjuntos convexos]] Indique verdadero (V) o falso (F):£I. Un segmento menos uno de los extremos es un conjunto convexo.£II. Una región triangular, menos una altura, es siempre un conjunto no convexo.£III. Algún ángulo es un conjunto convexo.£A)VVVæB)VFFæC)FFF£D)VFVææE)VVF£ [[clave=-]] [[Estado=sin_revisar]]
\end{enumerate}
